\begin{abstract}
Tokenization serves as a fixed interface between raw symbolic inputs and neural encoders, yet most existing tokenizers are optimized for large language models and natural language distributions.
When used as drop-in components in structured and semi-structured prediction pipelines, such tokenizers induce uncontrolled compression that distorts role-dependent semantics and wastes token budget under tight compute constraints.
We formalize tokenization as a constrained interface design problem and introduce \emph{directed semantic distortion}, a prefix-aligned measure of irrecoverable task-relevant information loss induced by the interface.
Based on this formulation, we propose \textbf{Controlled Interface Tokenization (CIT)}, which enforces a versioned interface contract with role-explicit boundaries and typed-value integrity, and constructs task-specialized vocabularies by maximizing compression efficiency within the contract under deterministic execution.
Directed semantic distortion is used to evaluate and calibrate semantic fidelity across tokenizers and budgets, yielding empirical rate--distortion frontiers.
Across synthetic and public structured classification tasks, CIT achieves superior rate--distortion trade-offs, higher accuracy under matched token budgets, and improved robustness to semantics-preserving re-serialization and prefix-limited prediction.
\end{abstract}
